\chapter{实验结果}\label{chap:results}

本文所有实验均在闭环仿真中进行。测试时，策略根据当前观测生成末端执行器任务空间参考，下层控制器执行后环境返回下一步观测；若回合达到最大步数仍未完成抓取和运输目标，则判定为失败。每个评估配置运行200个独立回合。BC、BC-RNN、ACT和Diffusion四类模仿策略在共同真值完整场景上比较；Diffusion和MPC专家进一步覆盖真值运动分解矩阵；含噪观测矩阵则覆盖四类模仿策略并与MPC专家对照。成功率、Wilson区间、最终距离、跟踪误差、阈值内比例和平滑度分别按式~\eqref{eq:success-rate}--\eqref{eq:smoothness-metric}计算。下面从总体比较、真值观测、含噪观测、过程指标和失败模式五个角度分析结果。

\section{模仿学习策略总体比较}

表~\ref{tab:imitation-baselines}给出了四类模仿学习策略在真值完整场景上的200回合结果。四种方法使用同一目标中心化数据表示、同一任务空间动作块接口和同一下层MPC执行器，因此该表主要反映不同学习模型结构对闭环任务的影响。

\begin{table}[!htbp]
    \caption{真值完整场景上的模仿学习策略比较。}
    \label{tab:imitation-baselines}
    \centering
    \small
    \begin{tabular}{lcccc}
        \hline
        策略 & 成功率 & 95\%区间 & 最终距离mean/median/p95 & 平滑度RMS \\
        \hline
        BC & 84.0\% (168/200) & [78.3\%, 88.4\%] & 0.1281/0.0240/0.8777 & 5.5863 \\
        BC-RNN & 81.0\% (162/200) & [75.0\%, 85.8\%] & 0.1354/0.0246/0.8877 & 5.8179 \\
        ACT & 81.5\% (163/200) & [75.5\%, 86.3\%] & 0.1578/0.0239/0.8860 & 5.7308 \\
        Diffusion & 86.5\% (173/200) & [81.1\%, 90.6\%] & 0.1201/0.0227/0.8872 & 5.4299 \\
        \hline
    \end{tabular}
\end{table}

结果显示，四种模仿策略都能在完整动态场景中形成有效闭环控制。Diffusion达到86.5\%成功率，是该表中最高的学习策略；BC达到84.0\%，仅低2.5个百分点，说明在目标中心化表示和下层MPC执行器存在时，简单监督行为克隆已经是强方法。ACT和BC-RNN分别达到81.5\%和81.0\%，与Diffusion差距为5.0和5.5个百分点。四类方法的最终距离中位数、跟踪误差和轨迹平滑度相近，说明当前任务中的主要差异更多来自失败回合比例和抓取时机，而不是成功回合附近的小幅跟踪误差。

\begin{figure}[!htbp]
    \centering
    \includegraphics[width=0.95\textwidth]{fig_imitation_baselines}
    \caption{真值完整场景中四类模仿学习策略的闭环性能比较。}
    \label{fig:imitation-baselines}
\end{figure}

\section{真值观测评估}

表~\ref{tab:gt-results}给出了真值观测矩阵的200回合结果。基于表~\ref{tab:imitation-baselines}中的共同基准表现，本文以Diffusion作为学习策略代表，并与MPC专家参考进行比较。完整运动场景中，Diffusion成功173次，成功率为86.5\%；MPC专家成功190次，成功率为95.0\%。这表明在无观测噪声的主任务上，显式阶段逻辑和MPC执行仍具有更高上限，同时也为表~\ref{tab:imitation-baselines}中的各类模仿策略提供了专家参考。

\begin{table}[!htbp]
    \caption{真值观测矩阵成功率。}
    \label{tab:gt-results}
    \centering
    \scriptsize
    \begin{tabular}{lccc}
        \hline
        运动场景 & 扩散策略 & MPC专家 & 差值 \\
        \hline
        Full & 86.5\% (173/200) & 95.0\% (190/200) & +8.5 \\
        X translation only & 92.5\% (185/200) & 96.5\% (193/200) & +4.0 \\
        Y translation only & 99.0\% (198/200) & 100.0\% (200/200) & +1.0 \\
        XY translation only & 93.0\% (186/200) & 95.5\% (191/200) & +2.5 \\
        Rotation only & 100.0\% (200/200) & 100.0\% (200/200) & 0.0 \\
        X + rotation & 92.5\% (185/200) & 94.5\% (189/200) & +2.0 \\
        Y + rotation & 94.5\% (189/200) & 99.0\% (198/200) & +4.5 \\
        \hline
    \end{tabular}
\end{table}

分解场景呈现出更细的结构。纯旋转场景中两类策略都达到100.0\%成功率，说明在固定中心附近处理目标自旋并不是当前任务的主要瓶颈。单轴平移和双轴平移场景中，Diffusion与MPC专家差距为1.0至4.0个百分点；组合场景中，$x$向平移加旋转差距为2.0个百分点，$y$向平移加旋转差距为4.5个百分点。完整场景差距扩大到8.5个百分点，说明当前动作块模仿策略已经学习到多数局部对齐行为，但在多自由度耦合运动下仍不如显式阶段控制稳定。

\begin{figure}[!htbp]
    \centering
    \includegraphics[width=0.95\textwidth]{fig_gt_motion_matrix}
    \caption{真值观测下运动分解场景的成功率比较。}
    \label{fig:gt-motion-matrix}
\end{figure}

\section{含噪观测评估}

表~\ref{tab:noisy-results}给出了含噪观测矩阵结果。该矩阵将四类模仿学习策略与MPC专家置于相同扰动条件下比较，用于判断噪声鲁棒性来自具体生成式模型，还是来自目标中心化动作块接口和历史窗口本身。

\begin{table}[!htbp]
    \caption{含噪观测矩阵成功率。}
    \label{tab:noisy-results}
    \centering
    \scriptsize
    \begin{tabular}{lccccc}
        \hline
        噪声条件 & BC & BC-RNN & ACT & Diffusion & MPC专家 \\
        \hline
        Full mild & 54.0\% & 65.0\% & 70.0\% & 70.0\% & 84.0\% \\
        Full moderate & 33.0\% & 51.0\% & 56.0\% & 58.0\% & 2.5\% \\
        Full strong & 15.0\% & 27.5\% & 32.5\% & 25.0\% & 0.0\% \\
        Gaussian only, moderate & 58.5\% & 62.5\% & 69.0\% & 67.5\% & 2.5\% \\
        Drift only, moderate & 70.5\% & 71.0\% & 73.0\% & 78.5\% & 93.5\% \\
        Outlier only, moderate & 64.5\% & 82.0\% & 84.5\% & 75.0\% & 95.0\% \\
        Dropout + latency, moderate & 83.5\% & 83.0\% & 88.0\% & 85.0\% & 95.5\% \\
        \hline
    \end{tabular}
\end{table}

结果中最突出的现象是高斯噪声下模型控制与学习策略的反转。Gaussian-only moderate场景中，MPC专家几乎完全失效，仅为2.5\%；四类模仿策略均显著更高，BC、BC-RNN、ACT和Diffusion分别达到58.5\%、62.5\%、69.0\%和67.5\%。完整中度噪声和强噪声下也出现类似趋势：MPC专家分别为2.5\%和0.0\%，而学习策略仍保持15.0\%至58.0\%的成功率。这表明，当观测中存在连续高频抖动时，直接使用当前观测的确定性控制器会把目标位姿噪声传递到参考构造和MPC目标中；历史窗口、目标中心化表示和动作块输出则相当于在学习到的动作流形上生成较平滑的参考。ACT的动作分块与时间集成\citep{zhao2023learning}以及Diffusion Policy的条件去噪动作生成\citep{chi2023diffusion}为这一现象提供了相应解释。

不同学习策略在含噪矩阵中的排序也更清楚。ACT在七个含噪配置中均是最强或并列最强的非扩散基线，并在强组合噪声、高斯扰动、离群点以及丢帧延迟条件下超过Diffusion。Diffusion在完整中度噪声和漂移扰动下略优于ACT，在轻度完整噪声下二者相同。BC-RNN通常优于BC，特别是在完整组合噪声和离群点扰动下更明显。由此可见，噪声鲁棒性不能简单归因于扩散生成机制；在本文任务中，目标中心化动作块表示、历史条件编码以及模型结构共同决定了含噪闭环表现。

\begin{figure}[!htbp]
    \centering
    \includegraphics[width=0.95\textwidth]{fig_noisy_robustness}
    \caption{含噪观测条件下学习策略与MPC专家的鲁棒性比较。}
    \label{fig:noisy-robustness}
\end{figure}

\section{补充过程指标}

为补充成功率矩阵之外的闭环行为描述，本文进一步统计最终距离、跟踪误差、平滑度和成功步数等过程指标。这里选取Diffusion的代表性配置，结果见表~\ref{tab:aux-metrics}。

\begin{table}[!htbp]
    \caption{代表性闭环过程指标。}
    \label{tab:aux-metrics}
    \centering
    \scriptsize
    \begin{tabular}{lcccc}
        \hline
        配置 & 最终距离mean/median/p95 & 跟踪误差mean & 平滑度RMS mean & 成功步数mean \\
        \hline
        GT Full & 0.1201/0.0227/0.8872 & 0.1075 & 5.4299 & 192.1 \\
        GT Rotation only & 0.0181/0.0179/0.0279 & 0.1089 & 5.2045 & 189.1 \\
        GT X + rotation & 0.0780/0.0202/0.7015 & 0.1101 & 5.5632 & 190.2 \\
        Noisy full mild & 0.2611/0.0267/0.9354 & 0.1030 & 22.3354 & 194.3 \\
        Noisy full moderate & 0.3699/0.0314/0.9886 & 0.0939 & 32.6923 & 199.1 \\
        Noisy full strong & 0.6243/0.8186/1.0201 & 0.0762 & 43.2697 & 218.0 \\
        Noisy Gaussian moderate & 0.2841/0.0271/0.9270 & 0.1059 & 32.8374 & 193.4 \\
        Noisy dropout + latency & 0.1428/0.0227/0.8907 & 0.1059 & 5.2964 & 189.9 \\
        \hline
    \end{tabular}
\end{table}

这些指标给出三个补充观察。第一，真值完整场景中最终距离均值为0.1201 m，而中位数仅为0.0227 m，说明失败回合造成了明显长尾。第二，最终距离的均值通常远大于中位数，尤其在完整含噪场景中95百分位接近或超过0.9 m，说明失败回合往往是目标逃逸或抓取相位错过后造成的大偏差，而不是终端附近的小误差。第三，完整噪声从轻度到强度时，平滑度RMS均值从22.3354增加到43.2697，成功步数均值也延长到218.0步，表明强噪声不仅降低成功率，也使成功回合的闭环轨迹更加抖动且完成更慢。

\begin{figure}[!htbp]
    \centering
    \includegraphics[width=0.95\textwidth]{fig_process_metrics}
    \caption{代表性场景中的闭环过程指标分布。}
    \label{fig:process-metrics}
\end{figure}

\section{失败模式分析}

真值结果说明，四类模仿学习策略已经能够在动态抓取任务上执行完整抓取流程，但与MPC专家参考相比仍存在性能差距。一个可能原因是专家状态机显式利用几何规则和阶段条件，而学习策略需要从固定窗口中隐式学习抓取时机。动作块生成有助于平滑控制，但如果闭爪事件或相位判断偏移，后续MPC执行器也难以完全修正。BC与Diffusion差距较小，提示当前目标中心化表示本身贡献很大；ACT和BC-RNN没有明显超过BC，说明更复杂序列模型在该数据规模和任务设置下还需要更细的训练或正则化设计。

含噪结果说明，模型控制参考和学习策略的弱点不同。MPC专家对高斯抖动极端敏感，可能因为目标位姿噪声被直接输入参考构造和求解过程，导致末端参考频繁变化。学习策略在此类噪声下更稳，但优势不是无条件的：在轻度噪声、漂移、离群点和丢帧延迟场景中，显式阶段逻辑和几何规则仍使MPC专家保持最高成功率。该结果也符合离线模仿学习关于分布偏移的基本分析\citep{ross2010efficient,ross2011reduction}：当慢漂移或延迟长期污染历史窗口时，策略无法通过在线专家反馈修正被偏置的上下文。ACT在多数含噪学习策略比较中表现最强，说明序列建模和动作块接口本身对噪声鲁棒性具有重要作用；Diffusion的优势主要体现在完整中度噪声和漂移扰动下。

\begin{figure}[!htbp]
    \centering
    \includegraphics[width=0.92\textwidth]{fig_failure_modes}
    \caption{动态抓取中的代表性成功与失败模式。}
    \label{fig:failure-modes}
\end{figure}

当前结果仍有三点局限。第一，所有结论来自仿真环境，尚未验证真实传感、接触和执行器误差下的迁移性能。第二，本文的时间指标主要反映仿真闭环执行过程，不能直接等价为真实硬件实时控制能力。第三，真值运动分解矩阵仍以Diffusion作为学习策略代表，尚未覆盖BC、BC-RNN和ACT，因此不同模仿方法在运动因素分解下的相对表现仍需进一步实验支持。
